\begin{dashboardbox}[一页结论：研究完成的边界，而非配置授权]
\begin{tabularx}{\textwidth}{L{3.3cm}X}
\textbf{覆盖对象} & 3 只 ETF（511010、510880、512890）+ 7 只公募 REIT（508069、508011、508068、508077、508055、508080、508096）；另有“受保障流动性”与“防御分散器”两项未命名类别。\\
\textbf{R0 证据状态} & 已归档产品法律文件、2025 年年报、2026Q2 报告、项目/持仓穿透和上交所 2026-08-12 同日快照；季度财务均明确标为\textbf{未经审计}。\\
\textbf{当前投资动作} & \qzero。\textbf{不新增、不替换、不配置任何 ETF 或 REIT}；并非因为某一项已被否定，而是每一项仍缺执行、当前估值、账户/税务或组合回放所必需字段。\\
\textbf{策略目标口径} & 5\% 是待检验的总回报研究目标，不是年化底线、预测或概率承诺；3\% 仅可指滚动 12 个月、税后费后、实际结算且可复投的净现金，不包括卖出本金、未实现收益、外部入金或借款。\\
\textbf{发布状态} & \textbf{研究框架通过；实施未获批准。} 本备忘录允许继续观察和补证，不允许把历史分派率、标签或单日成交额转化为买入理由。\\
\end{tabularx}
\end{dashboardbox}

\section{本次纠正了什么}

此前的策略框架定义了可能的资产类别与条件性参考权重，但它本身不能等同于底层标的尽调。本案把所有被点名的交易产品逐一放入同一审计框架：\textbf{产品是什么、实际拥有或收费什么、历史现金/收益来自哪里、估值能否成立、主要风险是什么、还缺什么以及能否改变 q。} 结果是十个产品均有可追溯产品卡；两项未命名类别被明确标为“不可研究、不可准入”，不再以类别名称掩盖缺口。

\begin{houseviewbox}[本机构观点]
收益标签不是风险分类。511010 是集中于五年期国债曲线的利率暴露；510880 的已披露 Q2 行业权重集中于采矿与金融；512890 的金融行业权重超过一半；7 只 REIT 则分别承受交通、北京社区商业、租赁住房、亦庄产业园和能源现金流驱动。它们不能仅因历史分派或“低波动/红利”名称被拼接成确定的 5\%/3\% 产品组合。
\end{houseviewbox}

\section{全产品决策矩阵}

\begin{exhibitbox}[表 1：已命名产品逐项处置（截至 2026-08-12）]
\scriptsize
\begin{tabularx}{\exhibitboxwidth}{L{1.35cm}L{2.7cm}L{3.2cm}X L{2.0cm}}
\toprule
\textbf{代码} & \textbf{产品} & \textbf{经济暴露} & \textbf{已完成的核心核验} & \textbf{当前状态} \\
\midrule
511010 & 国泰上证5年期国债ETF & 五年期国债、利率曲线 & 法律条款、费率、年报/Q2 NAV、债券占比及市场快照；缺久期/DV01/同日 IOPV & \qzero \\
510880 & 华泰柏瑞上证红利ETF & 红利股票；采矿、金融权重较高 & 完全复制规则、年报/Q2、分派记录、行业/前十大持仓及快照；缺当前持仓与税后现金 & \qzero \\
512890 & 华泰柏瑞中证红利低波ETF & 红利低波股票；金融权重较高 & 完全复制规则、年报/Q2、持仓/行业与快照；近三年未实施利润分配 & \qzero \\
508069 & 华夏南京交通高速公路REIT & 单项收费公路、车流/收费 & 年报/Q2、通行费、OCF、分派、项目债务线索及快照；缺期限/估值/深度 & \qzero \\
508011 & 嘉实物美消费REIT & 北京四处社区商业、租户租金 & 年报/Q2、出租率/WALT/租户集中、OCF/分派与快照；缺租户信用和债务勾稽 & \qzero \\
508068 & 华夏北京保障房REIT & 北京六处公租房、政策租金 & 年报/Q2、出租/收缴、扩募、OCF/分派与快照；缺扩募后可比估值 & \qzero \\
508077 & 华夏华润有巢租赁住房REIT & 上海三处市场化长租公寓 & 年报/Q2、出租率、租金、短 WALT、扩募与快照；缺续租压力测试 & \qzero \\
508055 & 汇添富上海地产租赁住房REIT & 上海两处保租房 & 年报/Q2、出租率、收缴、分派与快照；首个报告年度不完整 & \qzero \\
508080 & 中金亦庄产业园REIT & 亦庄产业园、工业租户 & 年报/Q2、出租率、最大租户、利润--现金差异与快照；缺现金桥和租约/债务包 & \qzero \\
508096 & 中航京能国际能源REIT & 光伏及水电、发电/电价/应收 & 更名/扩募、年报/Q2、OCF/分派与快照；年报 NAV 多口径未勾稽 & \qzero \\
\bottomrule
\end{tabularx}
\sourcenote{产品宇宙、ETF/REIT 核验包与上交所快照：\sourceid{S04--S28}。所有“缺口”均记录于 source exhaustion log。}
\end{exhibitbox}

\begin{exhibitbox}[图 1：十只已命名产品的经济暴露与当前研究状态]
\centering
\begin{tikzpicture}[
  x=1cm,y=1cm,
  product/.style={rounded corners=2pt, draw=navy, fill=blue!4, minimum width=2.15cm, minimum height=0.62cm, align=center, font=\scriptsize},
  reit/.style={rounded corners=2pt, draw=riskamber!80!black, fill=yellow!10, minimum width=2.15cm, minimum height=0.62cm, align=center, font=\scriptsize},
  exposure/.style={rounded corners=2pt, draw=steelgrey, fill=panelgrey, minimum width=2.35cm, minimum height=0.62cm, align=center, font=\scriptsize},
  blocked/.style={rounded corners=2pt, draw=riskred, fill=red!6, minimum width=2.15cm, minimum height=0.62cm, align=center, font=\scriptsize\color{riskred}},
  arrow/.style={-{Stealth[length=2mm]}, draw=mutedgrey, line width=.35pt}
]
\node[product] (p1) at (0,3.10) {511010\\五年期国债};
\node[product] (p2) at (0,2.20) {510880\\红利权益};
\node[product] (p3) at (0,1.30) {512890\\红利低波权益};
\node[reit] (p4) at (0,0.40) {508069\\收费公路};
\node[reit] (p5) at (0,-0.50) {508011\\社区商业};
\node[reit] (p6) at (0,-1.40) {508068 / 508077 / 508055\\租赁住房风险簇};
\node[reit] (p7) at (0,-2.30) {508080\\产业园租户};
\node[reit] (p8) at (0,-3.20) {508096\\水光能源};
\node[exposure] (e1) at (4.35,3.10) {利率曲线};
\node[exposure] (e2) at (4.35,1.75) {权益因子、金融/周期集中};
\node[exposure] (e3) at (4.35,0.00) {项目运营、租约/车流/资源};
\node[blocked] (b) at (8.30,0.00) {全部 q=0\\不构成配置};
\draw[arrow] (p1.east) -- (e1.west);
\draw[arrow] (p2.east) -- (e2.west);
\draw[arrow] (p3.east) -- (e2.west);
\draw[arrow] (p4.east) -- (e3.west);
\draw[arrow] (p5.east) -- (e3.west);
\draw[arrow] (p6.east) -- (e3.west);
\draw[arrow] (p7.east) -- (e3.west);
\draw[arrow] (p8.east) -- (e3.west);
\draw[arrow] (e1.east) -- (b.north west);
\draw[arrow] (e2.east) -- (b.west);
\draw[arrow] (e3.east) -- (b.south west);
\end{tikzpicture}
\sourcenote{穿透与资料状态：\sourceid{S04--S28}。蓝色为 ETF、黄色为公募 REIT；这是定性经济暴露图，不是相关性、收益预测或配置建议。}
\end{exhibitbox}

\section{未命名类别不是已尽调产品}

\begin{riskbox}[两个不可跳过的范围缺口]
\textbf{受保障流动性类别：}当前 CNY200,000 的防御性停泊表述没有指定存款、直接国债或货币市场产品，因而不能核验法律债权、存款保险合并口径、期限、提前支取、税费或账户资格。它只是停泊\emph{类别}，不可以被称为已选产品或已知收益率。\par
\textbf{防御分散器类别：}没有标的、非杠杆筛选规则、相关性、压力测试和流动性资料，不能假设其能够降低组合回撤，更不能为 3\% 现金目标贡献一个虚构收益率。两类均维持 q=0，先冻结有限候选清单，才有下一轮产品尽调。
\end{riskbox}

\section{从观察到可讨论配置：必须同时满足的八个门槛}

\begin{exhibitbox}[图 2：从产品研究到可讨论配置的八项串联门槛]
\centering
\begin{tikzpicture}[
  gate/.style={rounded corners=2pt, draw=navy, fill=panelgrey, minimum width=2.18cm, minimum height=0.82cm, align=center, font=\scriptsize},
  finalgate/.style={rounded corners=2pt, draw=riskred, fill=red!6, minimum width=2.18cm, minimum height=0.82cm, align=center, font=\scriptsize\color{riskred}},
  arrow/.style={-{Stealth[length=2mm]}, draw=mutedgrey, line width=.4pt}
]
\node[gate] (g1) at (0,1.20) {1\\法律身份与\\资产清单};
\node[gate] (g2) at (2.75,1.20) {2\\同日市场价与\\NAV / IOPV};
\node[gate] (g3) at (5.50,1.20) {3\\点差、深度与\\20 日容量};
\node[gate] (g4) at (8.25,1.20) {4\\账户、税务、\\结算与适配};
\node[gate] (g5) at (1.38,-0.40) {5\\久期/持仓或\\项目运营与债务};
\node[gate] (g6) at (4.13,-0.40) {6\\共同风险、\\地域与管理人压力};
\node[gate] (g7) at (6.88,-0.40) {7\\历史 TWRR 与\\已结算净现金回放};
\node[finalgate] (g8) at (9.63,-0.40) {8\\独立风险预算与\\组合层批准};
\draw[arrow] (g1.east) -- (g2.west);
\draw[arrow] (g2.east) -- (g3.west);
\draw[arrow] (g3.east) -- (g4.west);
\draw[arrow] (g4.south) |- (g5.east);
\draw[arrow] (g5.east) -- (g6.west);
\draw[arrow] (g6.east) -- (g7.west);
\draw[arrow] (g7.east) -- (g8.west);
\end{tikzpicture}
\sourcenote{门禁来自本案产品宇宙、回报模型、风险框架和资料穷尽日志。八项必须同时满足；图中顺序仅表达依赖关系，不表示任何产品将获准配置。}
\end{exhibitbox}

\begin{exhibitbox}[图 3：历史资料、市场快照与当前实施资料之间的断点]
\centering
\begin{tikzpicture}[
  x=1cm,y=1cm,
  observed/.style={rounded corners=2pt, draw=navy, fill=blue!4, minimum width=2.55cm, minimum height=1.02cm, align=center, font=\scriptsize},
  missing/.style={rounded corners=2pt, draw=riskred, fill=red!6, minimum width=2.70cm, minimum height=1.02cm, align=center, font=\scriptsize\color{riskred}},
  arrow/.style={-{Stealth[length=2mm]}, draw=mutedgrey, line width=.4pt},
  note/.style={font=\scriptsize\color{steelgrey}, align=center}
]
\node[observed] (annual) at (-5.05,0) {2025-12-31\\经审计年报\\历史 NAV / 现金};
\node[observed] (q2) at (-1.68,0) {2026-06-30\\未经审计 Q2\\持仓 / 项目运营};
\node[observed] (market) at (1.68,0) {2026-08-12 16:29\\上交所快照\\价格 / 成交额};
\node[missing] (current) at (5.15,0) {当前仍缺\\同日 NAV / IOPV、点差\\深度、容量与账户资料};
\draw[arrow] (annual.east) -- (q2.west);
\draw[arrow] (q2.east) -- (market.west);
\draw[arrow] (market.east) -- (current.west);
\node[note] at (0,-1.15) {日期错配与缺失字段意味着：不能将历史 NAV 诊断升级为当前估值，\\也不能将单日成交额升级为可执行容量。};
\end{tikzpicture}
\sourcenote{年报/Q2 与交易所快照：\sourceid{S04--S27}。所有十项的文件级日期、相对路径与 SHA-256 见 data/legal\_document\_manifest\_20260812.json；红色节点为实施阻断项。}
\end{exhibitbox}

未满足任何一项时，研究降级为观察；不以高收益替代、杠杆、降低信用质量或出售本金来“修复”目标。
